\subsection{按信息需求构建组合评价补充}
\label{sec:composition-data}

新来源基础题集侧重参数读取与条件保持。为进一步评价连接、筛选和判断等操作，本文在相同四个来源家族和156条读法目录上构建\CompositionQuestionCount{}道组合评价题，每个家族三题。编题前登记数量上限、来源页面、准入条件及运行预算。问题围绕工作模式比较、配置与量值对应、需求阈值筛选、环境范围和组合支持判断组织，并明确区分来源给出的规格与查询设定的工作点。

准入依据是信息需求：每题至少使用两个有独立作用的来源事实，且需要实际配对、筛选或判断。检查同时考察必要信息及约束被省略时，答案内容或选择结论如何变化；操作数量作为补充描述。对同一状态下的功耗与电流等相关表示，按工作状态归为事实组。该处理保留资料中的对应关系，使事实计量与实际信息依赖相符。

两名作者分别依据原件编写题目与预期答案，在程序编写前封存；一名非作者AI独立阅读全部八页资料，先对12题作答并封存，再比对作者参考与实际程序输出。12题的量值、条件及结论与原文一致。审阅中的一项必要性说明修订将相关量值按两个工作状态分组，题面、数值和程序保持原版本内容。全部修订与首次独立答案分别留档。

\begin{table}[htbp]
\centering\small
\caption{组合评价补充的来源分布与参考调用数}
\label{tab:composition-sources}
\CompositionSourceTable
\end{table}

组合题与原24道基础题分别冻结和统计，共享同一完整新来源目录。两部分具有相同来源家族和部分事实，来源覆盖按四个家族报告；训练与开发仍分别使用原576题和108题。组合题参考供执行核验及评价使用，域内恢复训练继续从训练题的真实模型失败中构造样本。
